Mobile GUI agents sometimes must type text they saw many screens earlier. We call such a step a \emph{delayed information need}. It is absent from AITZ (0 of 18,643 steps) and present in 31\% of sampled GUI-Odyssey trajectories, with a median gap of 5 steps between the source screen and the step that needs it. We test whether a \emph{screen index}, one LLM-written line per screen from which a reader picks a screen to re-read, beats keeping the full OCR history in context. On 138 verified cases it does not: full history localizes a source screen in 58.7\% of cases against 44.2\% for the index (\pval{0.008}) and types the correct text in 34.1\% against 20.3\% (\pval{0.002}). Training the writer from reader feedback does not close the gap. A rewritten line at the source screen reaches full-history level and no higher, and the 27-point gain of an oracle over five lines is redundancy that no single writer realizes. We also document two evaluation pitfalls, each larger than any writer effect we measured. Scoring only the earliest source screen is worth up to 15 points of localization. Keeping cases whose answer is already on the current screen turned a 17-point win for our own gating rule into a 23-point loss.
